\subsection{Теорема единственности}

\stepcounter{equation}
\begin{equation}\tag{1}\label{eq:uniq-problem}
  \begin{cases}
    \Box u \equiv u_{tt} - a^2 \Delta u = 0, & x \in \R^n,\ t > 0, \\[2pt]
    u|_{t=0} = 0, \quad u_t|_{t=0} = 0, \\[4pt]
    u \in C^2\big(\R^n \times (0,\infty)\big) \cap
         C^{2,1}\big(\R^n \times [0,\infty)\big) .
  \end{cases}
\end{equation}

\begin{warn}
На доске номер параграфа не был проставлен; в тетради он есть — \S4, так что
нумерация верна. Класс гладкости записан как
$C^{2,1}(\R^n \times [0,\infty))$ с пометкой «это условия для классического
решения»: по смыслу две производные по $x$ и одна по $t$, непрерывные вплоть
до границы $t=0$, чтобы начальные условия имели смысл.
\end{warn}

\begin{theorem}[единственность]
Единственное решение задачи~\eqref{eq:uniq-problem} — тождественный нуль:
$u \equiv 0$.
\end{theorem}

\begin{idea}
Это и есть теорема единственности для задачи Коши: если два решения имеют
одинаковые данные $\varphi, \psi$ и одинаковую правую часть $f$, то их разность
решает~\eqref{eq:uniq-problem} и, значит, равна нулю. Заметим, что $n$
произвольно — формул Даламбера и Кирхгофа здесь не требуется.
\end{idea}

\subsubsection*{Энергетическое тождество}

Продифференцируем плотность энергии:
\[
  \frac12 \dd{}{t}\Big(u_t^2 + a^2 |\nabla u|^2\Big)
  = u_t u_{tt} + a^2 \sum_{i=1}^{n} u_{x_i} u_{x_i t}
  = a^2 \sum_{i=1}^{n} \dd{}{x_i}\big(u_{x_i} u_t\big) ,
\]
то есть

\stepcounter{equation}
\begin{equation}\tag{2}\label{eq:energy-identity}
  \frac12 \dd{}{t}\Big(u_t^2 + a^2 |\nabla u|^2\Big)
  - a^2 \sum_{i=1}^{n} \dd{}{x_i}\big(u_{x_i} u_t\big) = 0 .
\end{equation}

\begin{idea}
Последнее равенство в цепочке — это само уравнение. Действительно,
\[
  a^2 \sum_i \dd{}{x_i}\big(u_{x_i} u_t\big)
  = a^2 \underbrace{\Delta u}_{=\,u_{tt}/a^2}\, u_t
    + a^2 \sum_i u_{x_i} u_{x_i t}
  = u_t u_{tt} + a^2 \sum_i u_{x_i} u_{x_i t} .
\]
Тождество~\eqref{eq:energy-identity} — закон сохранения энергии в
дифференциальной форме: $\partial_t(\text{плотность энергии}) +
\operatorname{div}(\text{поток энергии}) = 0$, где поток равен
$-a^2 u_t \nabla u$.
\end{idea}

\subsubsection*{Интегрирование по усечённому конусу}

Зафиксируем точку $(x_0,t_0)$ и возьмём усечённый конус
\[
  D = \big\{\, (x,t) : \ 0 \le t \le t_1, \ \ |x - x_0| \le a(t_0 - t) \,\big\},
  \qquad t_1 \in (0, t_0) .
\]
Его граница состоит из трёх кусков: нижнего основания $\Gamma_0$ (при $t=0$),
верхнего среза $\Gamma_1$ (при $t = t_1$) и боковой поверхности $\Gamma$.

\inputfig{figures/fig-energy-frustum}

Интегрируем~\eqref{eq:energy-identity} по $D$ и применяем формулу
Гаусса—Остроградского в $\R^{n+1}$; через $n = (n_1,\dots,n_n, n_t)$ обозначена
внешняя нормаль:
\begin{align*}
  0 &= \int\limits_D \left[\frac12 \dd{}{t}(\,\cdot\,)
       - a^2 \sum_{i=1}^{n} \dd{}{x_i}(\,\cdot\,)\right] dx\, dt \\[4pt]
    &= \int\limits_{\partial D} \left[\frac12 n_t
       \Big(u_t^2 + a^2 |\nabla u|^2\Big)
       - a^2 \sum_{i=1}^{n} n_i\, u_{x_i} u_t \right] ds
     = \int\limits_{\Gamma_0} \cdot
     + \int\limits_{\Gamma_1} \cdot
     + \int\limits_{\Gamma} \cdot
\end{align*}

Разберём три слагаемых по отдельности.

\paragraph{На $\Gamma_0$.} Нулевые начальные данные дают $u_t|_{t=0} = 0$ и
$\nabla u|_{t=0} = 0$, поэтому
\[
  \int\limits_{\Gamma_0} \cdot \ = 0 .
\]

\paragraph{На $\Gamma_1$.} Здесь нормаль равна $(0,\dots,0,1)$, значит
\[
  \int\limits_{\Gamma_1} \cdot
  = \frac12 \int\limits_{\Gamma_1} \Big(u_t^2 + a^2 |\nabla u|^2\Big) ds
  \ \ge\ 0 .
\]

\paragraph{На $\Gamma$.} Боковая поверхность характеристическая, и для её
нормали выполняется
\[
  n_t^2 = a^2 \sum_{i=1}^{n} n_i^2 .
\]
Умножим и разделим подынтегральное выражение на $n_t$ и воспользуемся этим
соотношением:
\begin{align*}
  \int\limits_{\Gamma} \cdot\ ds
  &= \int\limits_{\Gamma} \frac{1}{n_t}\left[
       \frac12 n_t^2 u_t^2
       + \frac{a^2}{2} \sum_{i=1}^{n} n_t^2 u_{x_i}^2
       - a^2 \sum_{i=1}^{n} n_t n_i\, u_{x_i} u_t
     \right] ds \\[4pt]
  &= \int\limits_{\Gamma} \frac{a^2}{2 n_t} \sum_{i=1}^{n}
     \big(n_i u_t - n_t u_{x_i}\big)^2 ds \ \ge\ 0 .
\end{align*}

\begin{idea}
Фокус ровно в одном месте: $\tfrac12 n_t^2 u_t^2 = \tfrac{a^2}{2}\sum_i n_i^2
u_t^2$ по характеристическому соотношению. После этого все три слагаемых
собираются в полный квадрат
\[
  \frac{a^2}{2 n_t} \sum_i \Big(n_i^2 u_t^2 - 2 n_t n_i u_{x_i} u_t
  + n_t^2 u_{x_i}^2\Big) .
\]
Знак обеспечен тем, что $n_t > 0$: конус раскрыт вниз, и внешняя нормаль к
боковой поверхности имеет положительную компоненту по времени.
\end{idea}

\subsubsection*{Завершение доказательства}

Сумма трёх неотрицательных слагаемых равна нулю, значит каждое из них нулевое. В
частности
\[
  \int\limits_{\Gamma_1} \Big(u_t^2 + a^2 |\nabla u|^2\Big) ds = 0
  \qquad \forall\, t_1 \in [0,t_0] .
\]
Подынтегральное выражение неотрицательно и непрерывно, поэтому внутри конуса
\[
  u_t = 0, \qquad u_{x_i} = 0, \quad i = \overline{1,n}
  \qquad\Longrightarrow\qquad u = \mathrm{const}
  \qquad\Longrightarrow\qquad u \equiv 0 ,
\]
и в частности $u(x_0,t_0) = 0$. Точка $(x_0,t_0)$ была произвольной.
\hfill$\blacksquare$

\begin{idea}
Доказательство попутно даёт то, что в \S1 и \S3 было вычитано из явных формул:
\textbf{конечность скорости распространения}. Всё рассуждение велось внутри
одного конуса, поэтому если данные равны нулю лишь на шаре
$|x - x_0| \le a t_0$, то $u(x_0,t_0) = 0$ всё равно — значение в вершине
определяется только конусом зависимости. Метод при этом не зависит от
размерности и переносится на уравнения с переменными коэффициентами, где
никаких явных формул нет.
\end{idea}
